\section{Getting it to work}

\label{sec:newsource}

This section describes how to configure your own cluster as a source that replicates into chango. Kueyen, the proof-of-concept source, is used as the worked example. Replace \texttt{<site>} with your cluster name and \texttt{<s3-endpoint>} with the hostname of your S3 endpoint.

\subsection{Prerequisites}

\begin{itemize}
  \item A Kubernetes cluster running Rook-Ceph, with the \texttt{rook-ceph} namespace and a healthy CephCluster.
  \item An nginx ingress controller and cert-manager with a \texttt{letsencrypt} ClusterIssuer, or another way to serve HTTPS.
  \item A DNS record for \texttt{<s3-endpoint>} (kueyen uses \texttt{s3.elkan.kueyen.dev.lsst.org}).
  \item HTTPS connectivity in both directions between your endpoint and \url{https://s3.chango.ls.lsst.org}.
  \item Access to the \texttt{rgw-realm-elkan} item in 1Password.
\end{itemize}

\subsection{Step 1: Realm credentials}

Every zone in the realm must use the same system keys. If your cluster runs External Secrets with a \texttt{onepassword} ClusterSecretStore, this is how kueyen does it:

{\small
\begin{verbatim}
---
apiVersion: external-secrets.io/v1
kind: ExternalSecret
metadata:
  name: elkan-keys
  namespace: rook-ceph
spec:
  secretStoreRef:
    kind: ClusterSecretStore
    name: onepassword
  target:
    name: elkan-keys
  data:
    - secretKey: access-key
      remoteRef:
        key: rgw-realm-elkan
        property: access-key
    - secretKey: secret-key
      remoteRef:
        key: rgw-realm-elkan
        property: secret-key
\end{verbatim}
}

\subsection{Step 2: Realm, zonegroup and zone}

{\small
\begin{verbatim}
---
apiVersion: ceph.rook.io/v1
kind: CephObjectRealm
metadata:
  name: elkan
  namespace: rook-ceph
spec:
  defaultRealm: false
---
apiVersion: ceph.rook.io/v1
kind: CephObjectZoneGroup
metadata:
  name: elkan
  namespace: rook-ceph
spec:
  realm: elkan
---
apiVersion: ceph.rook.io/v1
kind: CephObjectZone
metadata:
  name: elkan-<site>                # kueyen: elkan-kueyen
  namespace: rook-ceph
spec:
  zoneGroup: elkan
  customEndpoints:
    - https://<s3-endpoint>         # kueyen: s3.elkan.kueyen.dev.lsst.org
  metadataPool:
    failureDomain: host
    replicated:
      size: 3
    quotas:
      maxSize: 10Gi
  dataPool:
    failureDomain: host
    erasureCoded:
      dataChunks: 2
      codingChunks: 1
    quotas:
      maxSize: 1Ti
  preservePoolsOnDelete: true
\end{verbatim}
}

\subsection{Step 3: Object store}

{\small
\begin{verbatim}
---
apiVersion: ceph.rook.io/v1
kind: CephObjectStore
metadata:
  name: elkan
  namespace: rook-ceph
spec:
  preservePoolsOnDelete: true
  gateway:
    port: 80
    instances: 2
    resources:
      limits:
        cpu: "4"
        memory: 4Gi
      requests:
        cpu: "1"
        memory: 2Gi
  zone:
    name: elkan-<site>
\end{verbatim}
}

\subsection{Step 4: Ingress}

Keep the timeout and buffering annotations. Without them, large objects and long sync requests fail.

{\small
\begin{verbatim}
---
apiVersion: networking.k8s.io/v1
kind: Ingress
metadata:
  name: rook-ceph-rgw-elkan
  namespace: rook-ceph
  annotations:
    cert-manager.io/cluster-issuer: letsencrypt
    nginx.ingress.kubernetes.io/proxy-body-size: 1024m
    nginx.ingress.kubernetes.io/proxy-read-timeout: "3600"
    nginx.ingress.kubernetes.io/proxy-send-timeout: "3600"
    nginx.ingress.kubernetes.io/proxy-request-buffering: "off"
spec:
  ingressClassName: nginx
  tls:
    - hosts:
        - &host <s3-endpoint>
      secretName: rook-ceph-rgw-elkan-ingress-tls
  rules:
    - host: *host
      http:
        paths:
          - path: /
            pathType: Prefix
            backend:
              service:
                name: rook-ceph-rgw-elkan
                port:
                  number: 80
\end{verbatim}
}

\subsection{Step 5: StorageClass}

Buckets are provisioned on the source, never on chango but will appear on Chango due to the replication.

{\small
\begin{verbatim}
---
apiVersion: storage.k8s.io/v1
kind: StorageClass
metadata:
  name: elkan
provisioner: rook-ceph.ceph.rook.io/bucket
parameters:
  objectStoreName: elkan
  objectStoreNamespace: rook-ceph
reclaimPolicy: Retain
\end{verbatim}
}

\subsection{Step 6: Promote the zone to master}

The master role is not set by hand. It is set by adding two \texttt{radosgw-admin} commands to the cluster's \texttt{cephcli} ConfigMap in k8s-cookbook. Fleet deploys the ConfigMap with the rest of the \texttt{rook-ceph-conf} chart, and its script runs against the cluster, so the setting is version-controlled and reapplied automatically.

For kueyen, these lines were added to the end of the script in:
\path{fleet/lib/rook-ceph-conf/charts/kueyen/templates/cm-cephcli.yaml}:

{\small
\begin{verbatim}
---
apiVersion: v1
kind: ConfigMap
metadata:
  name: cephcli
  namespace: rook-ceph
data:
  script: |
    # ... existing cluster setup commands ...

    ceph osd pool set .mgr nodelete true
    ceph osd pool set .mgr nosizechange true

    # elkan multisite: kueyen is the master zone
    radosgw-admin zone modify --rgw-realm=elkan --rgw-zone=elkan-kueyen \
      --master --read-only=false --sync-from-all=false >/dev/null
    radosgw-admin period update --commit --rgw-realm=elkan >/dev/null
\end{verbatim}
}

The first command makes \texttt{elkan-kueyen} the writable master zone of the realm. The second commits a new period so the change reaches every zone, including chango. For a new source, add the same two lines to that cluster's \texttt{cm-cephcli.yaml}, with \texttt{elkan-<site>} as the zone name.

Chango's \texttt{cm-cephcli.yaml} holds the matching replica side:

{\small
\begin{verbatim}
    radosgw-admin zone modify --rgw-realm=elkan --rgw-zone=elkan \
      --read-only --sync-from-all=true >/dev/null
    radosgw-admin period update --commit --rgw-realm=elkan >/dev/null
\end{verbatim}
}

\subsection{Step 7: Point chango at your endpoint}

Open a pull request against k8s-cookbook the realm pull endpoint in
\texttt{fleet/lib/rook-ceph-conf/charts/chango/templates/cephobjectstore-elkan.yaml}:

{\small
\begin{verbatim}
apiVersion: ceph.rook.io/v1
kind: CephObjectRealm
metadata:
  name: elkan
  namespace: rook-ceph
spec:
  defaultRealm: true
  pull:
    endpoint: https://<s3-endpoint>
\end{verbatim}
}

\textbf{Note:} a zonegroup has exactly one master zone. If a master already exists (currently kueyen), a new source should not promote itself in Step 6. It should instead pull the realm from the existing master, by adding \texttt{pull.endpoint} to its own \texttt{CephObjectRealm}, and join as an additional writable zone. Chango already has \texttt{--sync-from-all=true}, so it replicates from every zone in the zonegroup. Coordinate with the team before adding a second source.

\subsection{Step 8: Test replication}

This demo writes to a dedicated \texttt{ittn-088} bucket on kueyen and shows it replicate to read-only chango.

\subsubsection{Create the demo bucket}

The bucket is a fixed-name owner plus an ObjectBucketClaim, both in \texttt{fleet/lib/rook-ceph-conf/charts/kueyen/templates/}.

{\small
\begin{verbatim}
# cephobjectstoreuser-ittn-088.yaml
---
apiVersion: ceph.rook.io/v1
kind: CephObjectStoreUser
metadata:
  name: ittn-088
  namespace: rook-ceph
spec:
  store: elkan
  clusterNamespace: rook-ceph
  quotas:
    maxBuckets: 1
\end{verbatim}
}

{\small
\begin{verbatim}
# obc-ittn-088.yaml
---
apiVersion: objectbucket.io/v1alpha1
kind: ObjectBucketClaim
metadata:
  name: &name ittn-088
  namespace: rook-ceph
spec:
  bucketName: *name
  storageClassName: elkan
  additionalConfig:
    bucketOwner: ittn-088
    bucketMaxSize: 1Gi
\end{verbatim}
}

\subsubsection{Set up aliases and credentials}

The aliases reach each cluster's toolbox and S3 endpoint, and the bucket owner's keys work on both because user metadata replicates.

{\small
\begin{verbatim}
alias T='kubectl --context kueyen -n rook-ceph exec deploy/rook-ceph-tools --'
alias C='kubectl --context chango -n rook-ceph exec deploy/rook-ceph-tools --'

S=rook-ceph-object-user-elkan-ittn-088
aws configure set aws_access_key_id "$(kubectl --context kueyen \
  -n rook-ceph get secret $S -o jsonpath='{.data.AccessKey}' \
  | base64 -d)" --profile ittn-088
aws configure set aws_secret_access_key "$(kubectl --context kueyen \
  -n rook-ceph get secret $S -o jsonpath='{.data.SecretKey}' \
  | base64 -d)" --profile ittn-088
aws configure set region elkan --profile ittn-088
aws configure set s3.addressing_style path --profile ittn-088

alias s3k='aws --profile ittn-088 \
  --endpoint-url https://s3.elkan.kueyen.dev.lsst.org'
alias s3c='aws --profile ittn-088 \
  --endpoint-url https://s3.chango.ls.lsst.org'
\end{verbatim}
}

\subsubsection{Check that the bucket exists on both clusters}

The bucket shows up on chango without being created there, because bucket metadata replicates.

{\small
\begin{verbatim}
kubectl --context kueyen -n rook-ceph get obc ittn-088 \
  -o jsonpath='{.status.phase}{"\n"}'                    # Bound
T radosgw-admin bucket list --rgw-realm=elkan | grep ittn-088
C radosgw-admin bucket list --rgw-realm=elkan | grep ittn-088
\end{verbatim}
}

\subsubsection{Demonstrate replication}

The file written on kueyen should be readable on chango, and the write to chango should fail with \texttt{AccessDenied}.

{\small
\begin{verbatim}
# 1. Write on kueyen (source)
echo "ITTN-088 replication demo $(date)" > ittn-088-demo.txt
s3k s3 cp ittn-088-demo.txt s3://ittn-088/

# 2. Read it on chango a few seconds later
s3c s3 ls s3://ittn-088/
s3c s3 cp s3://ittn-088/ittn-088-demo.txt -

# 3. Chango refuses writes (AccessDenied)
echo nope | s3c s3 cp - s3://ittn-088/nope.txt
\end{verbatim}
}

\subsubsection{Check from the admin side}

Chango should list \texttt{ittn-088-demo.txt} and report \texttt{data is caught up with source}, and the last command lists sync errors if it doesn't.

{\small
\begin{verbatim}
C radosgw-admin bucket list --rgw-realm=elkan --bucket=ittn-088 \
  | jq -r '.[] | "\(.meta.size)\t\(.name)"'
T radosgw-admin bucket sync status --rgw-realm=elkan --bucket=ittn-088
C radosgw-admin sync status --rgw-realm=elkan
C radosgw-admin sync error list --rgw-realm=elkan | head -30
\end{verbatim}
}
